\documentclass[a4paper, oneside]{article}
\usepackage{graphicx}
\usepackage{amsthm}
\usepackage{amsmath}
\usepackage{amssymb}
\usepackage[a4paper,
            bindingoffset=0.2in,
            left=2cm,
            right=2cm,
            top=2cm,
            bottom=2cm,
            footskip=.25in]{geometry}
\usepackage[italian]{babel}
\usepackage{pgfplots}
\usepackage{tabularx}
\usepackage{tikz}
\usepackage{wrapfig}
\usepackage{color}
\usepackage[d]{esvect}
\definecolor{page}{rgb}{0.129,0.157,0.212}
\pagecolor{page}
\color{white}
\graphicspath{ {./images/} }
\usetikzlibrary{shapes.geometric}
\usetikzlibrary{datavisualization}
\usetikzlibrary{datavisualization.formats.functions}
\usetikzlibrary{patterns}
\pgfplotsset{width=10cm,compat=1.18}

\title{Appunti di Lab II}
\author{Tommaso Miliani}
\date{02-03-26}

\begin{document}
\newtheoremstyle{theoremEnv}
                {}          % Space above
                {}          % Space below
                {\slshape}  % Body font
                {}          % Indent amount
                {\bfseries} % Head font
                {.}         % Punctuation after head
                {\newline}  % Space after theorem head
                {}          % Theorem head spec
\theoremstyle{theoremEnv}

\newtheorem{definition}{Definizione}[section]
\newtheorem{theorem}{Teorema}[section]
\newtheorem{lemma}{Proposizione}[section]
\newtheorem{observation}{Osservazione}[section]
\newtheorem{corollary}{Corollario}[theorem]
\newtheorem{example}{Esempio}[section]
\newtheorem{remark}{Enunciato}[section]

\maketitle

\section{Forza elettromotrice: definizione e misura}
Il \textbf{multimetro} è un oggetto che è in grado di misurare differenze di potenziali, 
correnti o resistenza di un dato circuito. Si deve capire ora come è definita la 
\textbf{forza elettromotrice} ($fem$). Il simbolo circuitale di questi oggetti è 
il seguente:

Un \textbf{generatore di forza elettromotrice} deve mantenere costante
la corrente che passa in un dato circuito (e dunque mantenere la 
differenza di potenziale costante in un certo circuito), mantenendo una 
zona più carica positivamente ed una più carica negativamente. Il generatore stesso
è un conduttore e dunque la corrente che fluisce nel circuito deve allo stesso modo fluire all'interno del
generatore. All'interno del generatore è presente dunque un campo elettrico $\vv{E_s}$ (conservativo), 
dato dal fatto che ci sono delle cariche in più o delle cariche in meno, dunque 
su tutto il circuito si ha che
\begin{gather*}
    \int_{\gamma} \vv{E} \cdot \vv{dl} = 0   
\end{gather*} 
Questo è quello che succede se si ha una lamina a facce piane e parallele
con distribuzione di carica $\sigma_+$ da una parte e $\sigma_-$ dall'altra,
allora il campo elettrico esterno è zero ma non vuol dire che non sia presente 
(infatti è l'integrale di linea che è zero ma il campo all'esterno è presente).
Esiste anche un campo $\vv{E_l}$:
\begin{gather*}
    \oint \vv{E_S} \cdot \vv{dl} + \vv{E_l} \cdot \vv{dl} = \int_{B}^{A} \vv{E_l} \cdot \vv{dl} \neq 0       
\end{gather*} 
Dove $\vv{E_l}$ è il campo elettrico interno al generatore di tensione. Il fatto che il 
campo interno al generatore sia non conservativo, non è preoccupante: queste forze ci devono essere per poter generare tensione
nel circuito. Si definisce allora questo campo non conservativo come 
\begin{align}
    \oint_{B}^{A}\vv{E_l} \cdot \vv{dl} = fem  
\end{align} 
La \textbf{fem} è definita come l'energia per unità di carica spesa 
dal generatore per mantenere la situazione di non equilibrio di un 
circuito all'interno del quale scorrono delle cariche.

Il fatto che la corrente scorra all'interno del generatore, potrà mai 
questo essere un oggetto che riesce sempre a portare le cariche 
senza resistenza? Ogni generatore ha una resistenza interna e la prima 
esperienza sta nella misura di una fem incognita e della sua resistenza 
interna. Sperimentalmente si osserva che, preso un generatore di fem e una resistenza
in un circuito aperto da $A$ a $B$. Se non circola corrente, è 
ovvio che viene mantenuta la differenza tra la parte carica positivamente e quella
negativamente, ed è anche ovvio che la fem è uguale alla differenza di potenziale tra 
il punto $A$ ed il punto $B$:
\begin{gather*}
    i = 0 \ \Longrightarrow \  fem = V_A - V_B
\end{gather*} 
Quello che si osserva in funzione della corrente è che, se si collegano 
i due punti $A$ e $B$ (dunque si chiude il circuito), essa decresce linearmente a causa di quella resistenza $r$ 
fino a che non decresce completamente. DUnque 
\begin{gather*}
    \Delta V = V_A - V_B = fem - ri
\end{gather*}
Dunque la corrente sale fino a che non arriva ad un valore massimo, oltre la quale il 
generatore non riesce più ad andare e inizia a scaldarsi fino a che ad un certo punto
non smette di mandare cariche (o esplode). Dunque ogni generatore inizialmente si comporta 
come un generatore di tensione, poi come un generatore di corrente. Il prodotto tra $I$ e $V$ è 
la potenza erogata dal generatore. 

\section{Misurare}
Il voltmetro, strumento che serve per misurare la tensione, ha una resistenza 
interna molto grande che viene messa in parallelo al circuito e misura la
differenza di potenziale tra il punto $A$ ed il punto $B$. Dunque, per la legge di Ohm
\begin{gather*}
    \Delta V = V_A - V_B = iR
\end{gather*}
Sapendo che il multimetro in dotazione ha $R = 11.1 M\Omega$, la corrente $i$ si determina a partire dalla 
relazione vista prima per la fem (legge sperimentale):
\begin{gather*}
    iR = fem - ri \ \Longrightarrow \ V_A - V_B = fem - ri
\end{gather*} 
Dunque
\begin{gather*}
    i = \frac{fem}{r + R}
\end{gather*}
TUttavia la differenza di potenziale è uguale alla caduta di potenziale di $iR$
perché è la resistenza più grossa, dunque
\begin{gather*}
    iR = \frac{fem R}{R + r} = V_A - V_B
\end{gather*}
Ossia
\begin{gather*}
    V_A - V_B  \approx fem \ \text{ se } R >> r
\end{gather*}
Le resistenze in serie sono equivalenti alla somma delle due resistenze, dunque
\begin{gather*}
    R_T = R_1 + R_2
\end{gather*}
Per un collegamento in parallelo invece, date due resistenze in serie
$R_1$ e $R_2$, la differenza di potenziale tra i due nodi prima del gruppo di resistenze
sarà uguale a 
\begin{gather*}
    \Delta V = i_1 R_1 = i_2 R_2
\end{gather*}
poiché tutti i punti devono essere allo stesso potenziale, dunque posso esprimere 
la somma in uscita come 
\begin{gather*}
    i_1 + i_2 = i \ \Longrightarrow \ i_1 = \frac{\Delta V}{R_1} , i_2 = \frac{\Delta V}{R_2}
\end{gather*}
Allora posso esprimere
\begin{gather*}
    \Delta V = i \frac{R_1 R_2}{R_1 + R_2}
\end{gather*}
Dunque
\begin{gather*}
    \frac{1}{R_T} = \frac{1}{R_1} + \frac{1}{R_2}
\end{gather*}

\subsection{Partitori di tensione}
La \textbf{legge del partitore di tensione} è una legge che permette di 
ottenere il rapporto tra la tensione in ingresso ed in uscita da un gruppo di resistenze.
Per un circuito in cui ci sono delle correnti in serie, si utilizza la seguente
\begin{gather*}
    V_A - V_C = \Delta V \qquad V_A - V_C = iR_1 + iR_2
\end{gather*}
Mentre
\begin{gather*}
    V_B - V_C = iR_2
\end{gather*}
Allora posso ricavare la corrente in uscita in funzione della corrente in entrata.
\begin{gather*}
    V_B - V_C = \Delta V_O = \frac{V_A - V_C}{R_1 + R_2} \cdot R_2
\end{gather*}
Dunque si ottiene la regola del partitore di tensioni
\begin{align}
    \Delta V_{o} = \Delta V_i \frac{R_2}{R_1 + R_2}
\end{align}


\section{Esperienza number One}
Gli obiettivi dell'esperienza sono i seguenti
\begin{itemize}
    \item Misurare $V_1$ con un voltmetro
    \item Misurare $V_2$ con due voltmetri.
\end{itemize}
Si vuole determinare, dato $V_1$ e $V_2$, 
la $fem$ del circuito. Si utilizza dunque la regola 
del partitore su di un circuito che presenta due resistenze in parallelo 
in serie con una resistenza (che è quella del generatore). Si possono considerare
le resistenze in parallelo come se fossero una sola di 
resistenza
\begin{gather*}
    R_T = \frac{R}{2}
\end{gather*}
Infatti
\begin{gather*}
    V_2 = \frac{\frac{R}{2}}{\frac{R}{2} + R_x} \epsilon_x
\end{gather*}
In quanto sono uguali. DUnque voglio determinare la resistenza interna del generatore 
non attraverso la definizione operativa ma attraverso 
\begin{gather*}
    \left\{\begin{array}{l}
        (R + 2R_x) V_2 = R \epsilon_x \\
        (R + R_x)V_1 = R\epsilon_x
    \end{array}\right.
\end{gather*}
Ottenendo allora 
\begin{gather*}
    R_x = R\left(\frac{V_2 - V_1}{V_1 - 2V_2}\right)
\end{gather*}
E
\begin{gather*}
    \epsilon_x = \frac{V_1(V_1 - 2V_2) + V_2V_1 - V_1^{2}}{V_1 - 2V_2}
\end{gather*}
Le incertezze degli strumetni sono
\begin{itemize}
    \item Multimetro FLuke 114: 0.5\% + 2 digit sulla tensione
    \item Multimetro Fluke 114 da $10k\Omega \sim 40 M\Omega$ è 1.5\% + 2 digit. 
\end{itemize}
Dove le percentuali sono gli errori relativi sulle misure (tutte) effettuate, 
che derivano dalla linearità dello strumento, mentre i "digit" dipendono 
dal digitalizzatore (ossia il convertitore dal segnale analogico al digitale), 
dunque sono le incertezze assolute sull'ultima cifra mostrata (di 2). 



\end{document}